% ===================================================================
%  نقشہ نمبر 15 — منڈی۔ — کھاݨ دیاں چیزاں۔
%  Traduction 2026 d'apres texto/io, controlee sur texto/fr, texto/en
%  et texto/hi. Consigne : voir texto/pnb/10-naqsha-01.tex.
%
%  LES FROMAGES ET LES LIQUEURS GARDENT LEUR NOM D'ORIGINE, transcrit
%  en shahmoukhi : ہالینڈ, کامانبیر, گروئیر, کِرش, کوراساؤ. Le
%  vocabulaire du livret est justement celui des choses etrangeres
%  qu'un marche apporte.
% ===================================================================

%%K t15-apar-1 apar
\VUsaut{4.0mm}
\VUtitre{15.0pt}{60}{نقشہ نمبر 15}
\VUsaut{3.0mm}

%%K t15-tit-1 sub
\VUpk{11.4pt}{40}{منڈی۔ --- کھاݨ دیاں چیزاں۔}
\VUsaut{3.0mm}

%%K t15-01-1 p
1. --- جیہڑے وپاری سانوں ساڈی لوڑ دا کھاݨا دیندے نیں، اوہناں وچوں بہتے
\VUgras{ڈھکی منڈی} دے \VUgras{ہال} \textsuperscript{(1)} دے تھلے جمع رہندے نیں، یا
کول دیاں گلیاں وچ۔

%%K t15-02-1 p
2. --- \VUgras{سُور دے مانس دا قصائی} \textsuperscript{(2)}، جیہڑا \VUgras{ہیم}
\textsuperscript{(3)}، \VUgras{لہو دی ساسج} \textsuperscript{(4)}، \VUgras{ساسج}
\textsuperscript{(5)}، \VUgras{آندراں دی ساسج} \textsuperscript{(6)} تے
\VUgras{چربی} \textsuperscript{(7)} ویچدا اے، اوہ \VUgras{مکھݨ تے پنیر ویچݨ والی}
\textsuperscript{(8)} کول بہندا اے، جیدا ٹھیلا لال چھِلّ والے
\VUgras{ہالینڈ دے پنیراں} \textsuperscript{(9)}، \VUgras{کامانبیر}
\textsuperscript{(10)}، \VUgras{گروئیر دے وڈے چکراں} \textsuperscript{(11)} تے تازے
\VUgras{مکھݨ دے پیڑیاں} \textsuperscript{(12)} نال سجیا اے۔

%%K t15-03-1 p
3. --- اوہدے سامݨے اک \VUgras{ترکاری ویچݨ والی} \textsuperscript{(13)} اپݨیاں
گاہکاں نوں سوہݨے \VUgras{ٹماٹر} \textsuperscript{(14)}، \VUgras{پھُل گوبھی}
\textsuperscript{(15)}، \VUgras{سلاد} \textsuperscript{(16)}، \VUgras{بند گوبھی}
\textsuperscript{(17)}، \VUgras{کدو} \textsuperscript{(19)}، \VUgras{خربوزے}
\textsuperscript{(18)} تے \VUgras{کھمبیاں} \textsuperscript{(20)} تیکر وکھاندی اے۔ پر
اک \VUgras{شراب والے} نے اپݨی \VUgras{ڈھوݨ دی گڈی} \textsuperscript{(21)}، جیہڑی
\VUgras{بیئر دے پیپیاں} \textsuperscript{(22)} نال لدی اے، ٹھیک اوہدے ٹھیلے دے سامݨے
کھڑی کر دِتی اے تے اوہدیاں گاہکاں نوں کول نئیں آوݨ دیندی۔ اک باغبان اوہنوں
\VUgras{ہاتھی چک} \textsuperscript{(23)} دی اک ٹوکری لیا رہیا اے۔

%%K t15-tit-2 sub
\VUsaut{4.0mm}
\VUpk{10.2pt}{40}{ویچݨا تے خریدݨا۔}
\VUsaut{3.0mm}

%%K t15-04-1 p
4. --- منڈی وچ وڈی چہل پہل اے، کیوں جے \VUgras{نیلامی} \textsuperscript{(24)} دا
ویلا آ گیا اے۔ جیہڑیاں \VUgras{گھر والیاں} \textsuperscript{(26)} ایتھے سودا کرن
آندیاں نیں، اوہ راہ وچ \VUgras{اوجھڑی ویچݨ والی} \textsuperscript{(25)} دے ٹھیلے دے
سامݨے رُک کے \VUgras{اوجھڑی} یا اک \VUgras{وچھے دا سر} خریددیاں نیں۔

%%K t15-05-1 p
5. --- اک موٹا \VUgras{رسوئیا} \textsuperscript{(27)} \VUgras{مچھی ویچݨ والی}
\textsuperscript{(28)} نال اک بہت سوہݨی \VUgras{ٹونا} دا مُل بھاء کر رہیا اے، جیہڑی
اپݨے من نال بھاء ودھاندی جاندی اے۔ اوہدے کول \VUgras{بام، سول، ٹراؤٹ} تے
\VUgras{کارپ} دی وڈی کھیپ آئی اے۔ اوہدا \VUgras{کاڈ} تے اوہدیاں \VUgras{سِپیاں}
بہت تازیاں نیں۔

%%K t15-tit-3 sub
\VUsaut{4.0mm}
\VUpk{10.2pt}{40}{سستا مال تے مہنگا۔}
\VUsaut{3.0mm}

%%K t15-06-1 p
6. --- اوہدے کول بیٹھی \VUgras{کُکڑیاں ویچݨ والی} \textsuperscript{(29)} بہت ایماندار
اے۔ جدوں اوہ کوئی \VUgras{ٹرکی}، کوئی \VUgras{ہنس}، کوئی \VUgras{بطخ} یا سدھارن
\VUgras{کُکڑی} ویچدی اے، تے واجب بھاء توں ودھ نئیں منگدی۔

%%K t15-07-1 p
7. --- چونک دے کونے اُتے، پہلی منزل اُتے، اک \VUgras{کپڑے دی دکان والی}
\textsuperscript{(30)} نے ہاݨ ای دکان کھولی اے، جِتھے خریدݨ والیاں نوں
\VUgras{میز پوش}، \VUgras{رومال، پرݨے} تے \VUgras{تولیاں} دی بہت ودھیا چوݨ ہمیشہ
ملدی اے۔ اوسے منزل اُتے اک \VUgras{چھُری بݨاؤݨ والے} \textsuperscript{(31)} نے اپݨی
کارگاہ لائی اے۔ اوہ منڈی دیاں زنانیاں دیاں \VUgras{چھُریاں} تیز کردا اے تے مالی لئی
سب توں کرڑے \VUgras{فولاد} توں \VUgras{چھانگݨ دیاں دِرانیاں} بݨاندا اے۔

%%K t15-tit-4 sub
\VUsaut{4.0mm}
\VUpk{10.2pt}{40}{کِریانے دی دکان۔}
\VUsaut{3.0mm}

%%K t15-08-1 p
8. --- اوہدی کارگاہ دے تھلے کجھ ویلا پہلاں اک وڈی کھاݨے دی دکان کھُلی اے۔ ایہہ ہُݨ
پرانے دناں والی اوہ سادی نِکی \VUgras{کِریانے دی دکان} \textsuperscript{(32)} نئیں
رہی، جیہڑی نِرا روز دا مال ویچدی سی — \VUgras{چا} \textsuperscript{(33)}،
\VUgras{چاکلیٹ} \textsuperscript{(34)}، \VUgras{ڈلی دی کھنڈ} \textsuperscript{(37)} تے
\VUgras{صابن} \textsuperscript{(38)}۔ ایتھے سب کجھ وِکدا اے :
\VUgras{کُرکُرے بسکٹ}، \VUgras{ڈبہ بند چیزاں} \textsuperscript{(35)}،
\VUgras{شراباں} \textsuperscript{(36)}، \VUgras{رم، برانڈی، کِرش، انیسیت} تے
\VUgras{کوراساؤ}۔ ایتھے شکار وی اوہنی ای سوکھ نال ملدا اے — \VUgras{سہا}
\textsuperscript{(39)}، \VUgras{چڑیاں} \textsuperscript{(40)}، \VUgras{تِتر}
\textsuperscript{(41)}، \VUgras{بٹیر} \textsuperscript{(42)}، \VUgras{جنگلی بطخ}
\textsuperscript{(43)} یا \VUgras{تِتر مور} \textsuperscript{(44)} — جِنّے نیمے
دیساں دے پھل، جِویں \VUgras{کیلے} \textsuperscript{(46)} تے \VUgras{انناس}۔

%%K t15-09-1 p
9. --- اک بہت مددگار کِریانے دا \VUgras{کارندہ} \textsuperscript{(45)} جیہڑا منگیا
جاوے اوہو دیݨ نوں تیار رہندا اے : \VUgras{مربہ} \textsuperscript{(47)}، کرونڈے دا یا
بہی دا، \VUgras{چربی} \textsuperscript{(48)}، \VUgras{کھیرے} \textsuperscript{(49)} یا
لُوݨی \VUgras{سارڈین} \textsuperscript{(50)}۔

%%K t15-09-2 p
ایہہ \VUgras{گاہکاں} \textsuperscript{(51)} لئی بہت سوکھ دی گل اے، جیہڑیاں سب توں
سدھارن مال دے کول ای سب توں ٹاواں اگیتا مال تے سب توں سوادی پھل لبھدیاں نیں : رسیلیاں
\VUgras{ناشپاتیاں} \textsuperscript{(52)}، \VUgras{آلو بخارے} \textsuperscript{(53)}،
\VUgras{انگور} \textsuperscript{(54)}، \VUgras{آڑو} \textsuperscript{(55)}،
\VUgras{بادام} \textsuperscript{(56)} تے \VUgras{اخروٹ} \textsuperscript{(57)}۔

%%K t15-tit-5 sub
\VUsaut{4.0mm}
\VUpk{10.2pt}{40}{گاہک۔}
\VUsaut{3.0mm}

%%K t15-10-1 p
10. --- \VUgras{چاول} \textsuperscript{(58)} تے \VUgras{مسور} \textsuperscript{(59)}
دھوئیں دِتی \VUgras{ہیرنگ} \textsuperscript{(60)} دے ڈبے کول رکھے نیں؛ \VUgras{آلو}
\textsuperscript{(61)} تے \VUgras{سیم} \textsuperscript{(63)} \VUgras{سیباں}
\textsuperscript{(62)}، \VUgras{انجیراں} \textsuperscript{(64)}،
\VUgras{سُکے آلو بخاریاں} \textsuperscript{(65)} تے \VUgras{پہاڑی باداماں}
\textsuperscript{(66)} کول بیٹھے نیں۔ ایس دکان وچ تازے \VUgras{آنڈے}
\textsuperscript{(67)} تے \VUgras{زیتون} \textsuperscript{(68)} وی ملدے نیں، سو
\VUgras{ٹوکریاں} \textsuperscript{(70)} تے \VUgras{جالی دے تھیلے} \textsuperscript{(73)}
بہت چھیتی بھر جاندے نیں۔

%%K t15-11-1 p
11. --- ایس ودھیا دکان دی \VUgras{کافی} \textsuperscript{(72)} نے
\VUgras{نوکراݨیاں} \textsuperscript{(69)} تے رسوئݨاں وچ چنگا خاصا ناں کمایا اے، کیوں
جے بُڈھا \VUgras{کِریانے والا} \textsuperscript{(71)} اوہنوں آپ بڑے دھیان نال بھُندا
اے۔

%%K t15-tit-6 sub
\VUsaut{4.0mm}
\VUpk{10.2pt}{40}{ویکھݨ دیاں چیزاں۔}
\VUsaut{3.0mm}

%%K t15-12-1 p
12. --- سارے لوک منڈی سودا کرن ای نئیں آندے۔ ہال تیکر جاندی اوس نِکی گلی دی سوہݨی
آواجائی وچ کئی طرحاں دے لوک وکھائی دیندے نیں۔ اک بھلے \VUgras{کمیلے دے مزدور}
\textsuperscript{(75)} کول، جیہڑا اپݨے اگے اک \VUgras{ٹھیلا} \textsuperscript{(74)}
دھکّ رہیا اے، ایہہ رہے دو چھٹی اُتے آئے سپاہی، جیہناں نوں اپݨی چمکدی وردی وکھاؤݨ دا
بڑا مان اے : \VUgras{ہوسار} \textsuperscript{(76)} اپݨے نیلے \VUgras{دولمان} تے
\VUgras{پراں والے شاکو} وچ، تے \VUgras{زُواو دا حوالدار} ڈھِلی نِکر وچ، سر اُتے
\VUgras{فیز} لائی۔

%%K t15-13-1 p
13. --- \VUgras{نانبائی} \textsuperscript{(80)}، جیہڑا \VUgras{نانبائی خانے}
\textsuperscript{(78)} دی دہلیز اُتے کھڑا اے، اوہنے ہُݨے ہُݨے اپݨی \VUgras{روٹی}
\textsuperscript{(79)} دا آٹا گُنھیا اے تے تندور توں اوہ \VUgras{کروسان}
\textsuperscript{(81)}، \VUgras{نِکیاں روٹیاں} \textsuperscript{(82)} تے
\VUgras{گول روٹیاں} \textsuperscript{(83)} کڈیاں نیں جیہڑیاں کھڑکی وچ رکھیاں نیں۔

%%K t15-14-1 p
14. --- نانبائی خانے دے اُتے اک ہوشیار \VUgras{سیوݨ والی} \textsuperscript{(84)}
رہندی اے، جیہڑی کئی \VUgras{کڑھائی کرن والیاں} \textsuperscript{(85)} رکھدی اے۔
اوہدے کول اک \VUgras{دھوبݨ تے استری کرن والی} \textsuperscript{(86)} رہندی اے، جیہڑی
\VUgras{کپڑے} \textsuperscript{(88)} دھو کے تے سُکا کے اوہناں نوں کلف لاندی اے تے
\VUgras{استری} \textsuperscript{(87)} نال دبدی اے۔

%%K t15-tit-7 sub
\VUsaut{4.0mm}
\VUpk{10.2pt}{40}{مانس، مچھی تے ترکاری۔}
\VUsaut{3.0mm}

%%K t15-15-1 p
15. --- ہیٹھلی منزل دی \VUgras{قصائی دی دکان} \textsuperscript{(89)} وچ ہر طرحاں دا
مانس وِکدا اے — \VUgras{بھیڈ دا} \textsuperscript{(90)}، \VUgras{گاں دا}
\textsuperscript{(94)} یا \VUgras{وچھے دا} \textsuperscript{(92)} — \VUgras{راناں،
چاپاں، ٹوٹے} تے \VUgras{پسلیاں} دے رُوپ وچ۔ \VUgras{قصائی دا چھوکرا}
\textsuperscript{(93)} ایہہ سارا مانس کٹدا اے تے \VUgras{گُدے دیاں ہڈیاں}
\textsuperscript{(96)} وکھ رکھدا اے۔ قصائی دے دروازے اُتے اک \VUgras{سِپیاں ویچݨ والی}
\textsuperscript{(97)} کھڑی اے، جیہڑی \VUgras{سِپیاں} \textsuperscript{(98)} ویچدی اے۔
چاہو تے اوہ اوہناں نوں کھول وی دیندی اے۔

%%K t15-16-1 p
16. --- ایہہ سارے وپاری لائسنس یا تھاں دا محصول دیندے نیں؛ ایسے لئی \VUgras{سپاہی}
\textsuperscript{(100)} اوہناں وچاریاں \VUgras{گھُم کے ترکاری ویچݨ والیاں}
\textsuperscript{(99)} دا بغیر ترس دے پچھا کردا اے، جیہڑیاں اپݨیاں ریڑھیاں اُتے
\VUgras{گاجراں، مُولیاں، شلغم، گندنا} یا \VUgras{شتاوری دیاں گٹھڑیاں، تازیاں مٹراں،
پالک، کھٹی بُوٹی، جل کُنبھی} تے \VUgras{دھنیا} ڈھو کے اوہناں نال مقابلہ کردیاں نیں۔

%%K t15-tit-8 sub
\VUsaut{4.0mm}
\VUpk{10.2pt}{40}{سپاہی توں بچو!}
\VUsaut{3.0mm}

%%K t15-17-1 p
17. --- \VUgras{تُھوم} \textsuperscript{(101)} تے \VUgras{گنڈے} ویچݨ والا مُنڈا خوب
جاݨدا اے پئی اوہنوں وی منڈی کول اپݨا مال ویچݨ دی اجازت نئیں۔ اوہ نس کھڑدا اے، ایس
خطرے نال پئی \VUgras{جھینگے مچھی} تے \VUgras{لابسٹر} \VUgras{ویچݨ والی}
\textsuperscript{(102)} دی \VUgras{کیکڑیاں}، \VUgras{کریفش} تے \VUgras{جھینگیاں} دی
ٹوکری اُلٹ دیوے۔

%%K t15-18-1 p
18. --- سارے لوک کاہلی وچ نیں تے وڈا رولا پا رہے نیں؛ پر ایس نال گھُمدے
\VUgras{شیشہ ساز} \textsuperscript{(103)} دی ہاک صاف سُݨائی دیوݨوں نئیں رُکدی، نہ اوس
\VUgras{کوئلے والے} \textsuperscript{(104)} دی، جیہنے ہُݨے ہُݨے اک
\VUgras{کوئلے دی بوری} \textsuperscript{(105)} اپڑاؤݨ نوں اپݨی گڈی تھوڑی دیر لئی
چھڈی اے۔
\VUsaut{6.0mm}
\VUfilet{22mm}
